%% 05-results.tex --- included by ../sfc.tex
\section{Results}
\label{sec:results}

Each subsection answers one research question of
Section~\ref{sec:rqs}. Protocol names and budgets are those of
Table~\ref{tab:budgets}; lower EPR is better throughout.

\subsection{RQ1: End-to-End Cost Under Each Compiler's Own Protocol}
\label{sec:rq1}

\emph{Question:} how do the three compilers compare when each uses its own
recommended protocol? \emph{Held fixed:} nothing --- each compiler supplies
its own layouts and pass schedule (Table~\ref{tab:budgets}).
\emph{Metric:} geometric-mean EPR over jointly completed circuits, plus
completion counts and complete compile time.

\label{sec:main_results}


\begin{table}[t]
\centering
\caption{Main comparison, own protocol: each compiler from its own
  initial layout. \TeleSABRE{} and \dSABRE{} report the best of three
  seeds; \PORTER{} routes one deterministic layout at each of three
  candidate lookahead horizons (Section~\ref{sec:horizon}), and
  \PORTER{}+R spends five routes on the backward-refinement schedule of
  Section~\ref{sec:sfc-refine} instead --- closer to \dSABRE{}'s nine to
  eighteen, though still short of them. Richer schedules that do reach
  that budget, and pass \dSABRE{}, are in Table~\ref{tab:protocols}.
  $\Delta_\mathrm{TS}$ is the EPR change of whichever \PORTER{} column is
  stronger on that row, against \TeleSABRE{} (\%); negative is better.
  The bolded \PORTER{} column is that stronger one --- refinement leads at
  25 and 64 qubits, horizon selection alone at 80. $^{\ast}$\TeleSABRE{} does not converge in
  three seeds; each gmean row covers the circuits it completes. Local
  SWAP counts, which follow the same ordering throughout, are tabulated
  in \appTables{}.}
\label{tab:main}
\setlength{\tabcolsep}{4pt}
\small
\resizebox{\columnwidth}{!}{%
\begin{tabular}{@{}ll rrrr r@{}}
\toprule
Suite & Circuit & \TeleSABRE{} & \dSABRE{} & \PORTER{} & \PORTER{}+R & $\Delta_\mathrm{TS}$ \\
\midrule
25q & ae          &        119 &    65 &    65 &    65 &   $-45.4$ \\
25q & ghz         &         13 &     4 &     4 &     4 &   $-69.2$ \\
25q & graphstate  &         21 &     6 &     7 &     4 &   $-81.0$ \\
25q & qft         &        138 &    81 &    74 &    74 &   $-46.4$ \\
25q & qnn         &        515 &   143 &   119 &   116 &   $-77.5$ \\
25q & random      &        636 &   404 &   401 &   383 &   $-39.8$ \\
\multicolumn{2}{@{}l}{\quad\textbf{gmean} (6)} & 106.6 & 44.0 & 43.1 & \textbf{38.8} & $\mathbf{-63.6}$ \\
\midrule
64q & ae          &        353 &   165 &   157 &   177 &   $-55.5$ \\
64q & ghz         &         16 &     6 &     6 &     6 &   $-62.5$ \\
64q & graphstate  &         62 &    15 &    20 &    16 &   $-74.2$ \\
64q & qft         &        312 &   179 &   262 &   262 &   $-16.0$ \\
64q & qnn         &        889 &   538 &   638 &   537 &   $-39.6$ \\
64q & random      &  $^{\ast}$ &   732 &   737 &   737 &       --- \\
64q & qpeexact    &        289 &   187 &   222 &   222 &   $-23.2$ \\
64q & qaoa        &        956 &   555 &   775 &   607 &   $-36.5$ \\
64q & multiplier  &       1601 &  1285 &  1076 &   763 &   $-52.3$ \\
\multicolumn{2}{@{}l}{\quad\textbf{gmean} (8)} & 284.5 & 144.6 & 166.2 & \textbf{149.2} & $\mathbf{-47.6}$ \\
\midrule
80q & ae          &  $^{\ast}$ &   332 &   299 &   490 &       --- \\
80q & ghz         &         26 &     7 &     7 &     7 &   $-73.1$ \\
80q & graphstate  &         87 &    20 &    28 &    24 &   $-72.4$ \\
80q & qft         &  $^{\ast}$ &   315 &   446 &   393 &       --- \\
80q & qnn         &  $^{\ast}$ &   791 &  1070 &  1181 &       --- \\
80q & random      &  $^{\ast}$ & 10551 & 11040 & 11040 &       --- \\
80q & qpeexact    &        784 &   347 &   291 &   526 &   $-62.9$ \\
80q & qaoa        &  $^{\ast}$ &   862 &  1122 &  1122 &       --- \\
80q & multiplier  &  $^{\ast}$ &  3793 &  2411 &  2699 &       --- \\
\multicolumn{2}{@{}l}{\quad\textbf{gmean} (3)} & 121.0 & 36.5 & \textbf{38.5} & 44.5 & $\mathbf{-68.2}$ \\
\bottomrule
\end{tabular}%
}
\end{table}



Table~\ref{tab:main} gives the per-circuit results. Against
\TeleSABRE{} the ordering is consistent across the suite. From one
deterministic layout \PORTER{} lowers geometric-mean EPR by
$59.6\%$, $41.6\%$ and $68.2\%$ at 25, 64 and 80 qubits, spending three
routes against \TeleSABRE{}'s three seeded attempts. Every circuit
\TeleSABRE{} completes improves, from $-16.0\%$ to $-81.0\%$.
Refinement adds to this at 25 and 64 qubits ($63.6\%$ and $47.6\%$) but
not at 80, where the selected route is already stronger ($68.2\%$
against $63.2\%$) --- the first place where choosing the horizon beats
spending four more routes. \PORTER{} and \dSABRE{} complete all
twenty-four instances; \TeleSABRE{} completes seventeen, failing on
64q~Random and six of the nine 80-qubit circuits, including the
$23{,}381$-CX Random.

Two details are worth recording. On all eight 64-qubit circuits
\TeleSABRE{} completes, our harness reproduces the EPR and SWAP counts
\dSABRE{} publishes for it~\cite{lidsabre} exactly, so the two
evaluations share a baseline measurement rather than only a baseline
tool. And GHZ, the chain circuit that defeats locally-greedy
schedulers, costs \TeleSABRE{} 16 pairs on 63 CX gates but \PORTER{} and
\dSABRE{} 6 each: \PORTER{}'s first-appearance visit order puts
successive chain neighbours in one core before routing begins.

Against \dSABRE{} the answer depends on the pass budget. At its own
three-route protocol \PORTER{} takes the geometric mean at 25 qubits
($43.1$ against $44.0$) and trails at 64 and 80 ($166.2$ against
$144.6$; $38.5$ against $36.5$, both on the \TeleSABRE{}-matched
basis), winning eight instances, losing twelve, tying four. Refinement
comes within $3.2\%$ at 64; the $\gamma$-search closes the 64-qubit
suite outright ($172.7$ against $173.1$ over all nine) and leaves
$1.08\times$ at 80. Its widest win is 64q~Multiplier, $763$ against
$1{,}285$ with refinement; its worst losses are $1.5\times$, on
80q~QPEexact ($526$ against $347$) and 64q~QFT ($262$ against $179$).
The first is a loss for refinement alone --- the selected route takes
that instance at $291$.

\subsection{RQ2: Routing Quality With the Layout Controlled}
\label{sec:shared_results}

\emph{Question:} with the initial layout held identical and the routing
budget matched, does either routing loop use fewer EPR pairs?
\emph{Held fixed:} the initial layout, supplied externally.
\emph{Metric:} geometric-mean EPR, and win/loss counts per instance.


The own-protocol comparison mixes two effects: the quality of each
compiler's initial layout and the quality of its routing loop.
Table~\ref{tab:shared} removes the first --- both \dSABRE{} and
\PORTER{} route, single-shot, from the identical layout \TeleSABRE{}
built for itself --- leaving a pure routing comparison.

\begin{table}[t]
\centering
\caption{Shared-layout comparison: \dSABRE{} and \PORTER{} both routed
  from the identical layout \TeleSABRE{} built for itself, so the only
  difference left is the routing loop. Geometric-mean EPR
  over the circuits \TeleSABRE{} completes and therefore provides a
  layout for --- all six at 25q, eight of nine at 64q, three of nine at
  80q. $\Delta_\mathrm{TS}$ columns are each compiler's change against
  \TeleSABRE{}'s own-protocol result. Per-circuit rows are in
  \appTables{}.}
\label{tab:shared}
\setlength{\tabcolsep}{4pt}
\small
\begin{tabular}{@{}l r rr rr@{}}
\toprule
 & & \multicolumn{2}{c}{\dSABRE{}} & \multicolumn{2}{c}{\PORTER{}} \\
\cmidrule(lr){3-4}\cmidrule(lr){5-6}
Suite & \TeleSABRE{} & EPR & $\Delta_\mathrm{TS}$ & $+$R & $\Delta_\mathrm{TS}$ \\
\midrule
25q (6) & 106.6 & 51.0 & $-52.2$ & \textbf{47.7} & $\mathbf{-55.2}$ \\
64q (8) & 284.5 & 189.2 & $-33.5$ & \textbf{157.2} & $\mathbf{-44.7}$ \\
80q (3) & 121.0 & 74.2 & $-38.7$ & \textbf{55.7} & $\mathbf{-53.9}$ \\
\textbf{all} (17) & 173.0 & 100.9 & $-41.7$ & \textbf{86.0} & $\mathbf{-50.3}$ \\
\bottomrule
\end{tabular}
\end{table}


Controlling the layout isolates the routing loop, and this is the
experiment \PORTER{}'s own protocol loses. At three routes each,
\dSABRE{} leads at 25 and 80 qubits and \PORTER{} at 64 ($-52.2\%$
against $-41.5\%$, $-38.7\%$ against $-33.1\%$, $-33.5\%$ against
$-36.5\%$), pooling to $100.9$ against $107.7$ over seventeen instances.
Spending three routes on lookahead horizons is therefore not what closes
the gap. Spending five on the forward--backward--forward schedule is:
\PORTER{}+R reaches $86.0$ against $100.9$ and leads at every scale.
That schedule, not the commitment rule, is what separates the two
compilers on fixed layouts, and Section~\ref{sec:ablation} confirms it
from the other side --- removing commitment moves EPR by $0.6\%$, in
\PORTER{}'s favour. (Table~\ref{tab:shared} reports the refined column;
the selected one is in \appTables{}.)

\textbf{Pooled fixed-layout grid.} A second source of externally
supplied layouts --- \texttt{pytket-dqc}'s hypergraph placements,
routed by both compilers on the 25q and 64q suites (\appPytket{};
\appOptin{} offers that partition to \PORTER{} as an extra layout
candidate instead) --- pools with Table~\ref{tab:shared} to give
forty-one instances on which \dSABRE{} and \PORTER{} routed the
\emph{identical} layout. Under \PORTER{}'s own protocol \dSABRE{} wins
twenty to sixteen with five ties, pooled geometric mean $134.2$ against
$139.0$; with the refinement schedule that matches its pass count
\PORTER{} wins twenty to fifteen with six ties, $126.6$ against
$134.2$. Routing quality, with the layout controlled and the pass
budget matched, does not separate these two compilers.

\textbf{Each compiler on the other's layout.} Closing the remaining
cell --- routing each system from its own and from its counterpart's
layout (\appTables{}) --- gives the same reading from the other side.
Down a column the compilers are close: on \dSABRE{}'s layout,
\dSABRE{} reaches $163.9$ and \PORTER{}+R $161.1$; on \PORTER{}'s,
$171.2$ and $171.4$. Across a row the layouts are not: \dSABRE{}'s
candidate set is the better start for \emph{both}, by $4.3\%$ for its
own router and $6.0\%$ for \PORTER{}'s. The best cell is \PORTER{}+R on
\dSABRE{}'s layout.

At these sizes, then, the two routing stages are worth about the same
given the same start, and what separates them end-to-end is that
\dSABRE{} searches up to six candidate layouts where \PORTER{} builds
one --- a deliberate trade that Section~\ref{sec:frontier} prices. It
does not hold at every size: from $150$ physical qubits upward the same
control becomes counter-productive and a different parameter binds
(Section~\ref{sec:scaleup}).

\subsection{RQ3: Mechanism Ablation}
\label{sec:ablation}

\emph{Question:} what does each mechanism contribute on its own?
\emph{Held fixed:} the initial layout, the lookahead horizon and every
other parameter; exactly one mechanism is disabled per column.
\emph{Metric:} EPR, completion, and --- for plan commitment, whose claim
is about cost --- compile time.


\begin{table}[t]
\centering
\caption{Mechanism ablation: \PORTER{} EPR with one design decision
  disabled at a time, all from the same temporal locality-greedy
  layout. Columns: ``No pk'' disables packets
  (\texttt{max\_packet\_size}=1), ``No rel.'' disables
  both relief tiers, ``Hop-wise'' replaces plan commitment with
  per-iteration re-scoring of the whole front, ``\TeleData{}'' removes
  the \TeleGate{} family; ``pk'' and ``tg'' are the packet and
  single-telegate counts of the full configuration. The last gmean row
  pairs Full against Hop-wise over all twelve rows, which both
  complete. $^{\dagger}$Fails to complete (deadlock
  or iteration-guard trip). $^{\ddagger}$Completes only with the additional
  relief tier described in the text, which the \TeleGate{} otherwise
  makes unnecessary. The first gmean row is over the fourteen instances the
  first three configurations share; the second is the \TeleData{}-only column
  paired against Full over the sixteen instances it completes.}
\label{tab:ablation}
\setlength{\tabcolsep}{3pt}
\small
\resizebox{\columnwidth}{!}{%
\begin{tabular}{@{}ll rrrrr rr@{}}
\toprule
Suite & Circuit & Full & No pk & No rel. & Hop-wise & \TeleData{} & pk & tg \\
\midrule
25q & ae         &    65 &    71 &    80 &    65 & 68$^{\ddagger}$ &   8 &    0 \\
25q & qft        &    74 &    74 &   140 &    74 &   74 &   0 &    0 \\
25q & qnn        &   119 &   122 &   212 &   128 &  119 &   0 &    0 \\
25q & random     &   401 &   412 &   395 &   399 & 413$^{\ddagger}$ &  43 &   93 \\
64q & ae         &   177 &   211 &   327 &   172 & 198$^{\ddagger}$ &  32 &    0 \\
64q & qft        &   262 &   372 &   382 &   336 & \textbf{216}$^{\ddagger}$ &  45 &    1 \\
64q & qnn        &   864 &   931 &   600 &   785 & \textbf{694}$^{\ddagger}$ & 240 &    5 \\
64q & random     &   737 &   843 &   789 &   805 & $^{\dagger}$ &  77 &  228 \\
80q & ae         &   490 &   688 & $^{\dagger}$ &   383 & \textbf{354}$^{\ddagger}$ & 264 &   15 \\
80q & qft        &   446 &   435 & $^{\dagger}$ &   415 & \textbf{307}$^{\ddagger}$ & 209 &    2 \\
80q & qnn        &  1517 &  1571 & $^{\dagger}$ &  1474 & \textbf{926}$^{\ddagger}$ & 448 &   30 \\
80q & random     & 11040 & 11628 & $^{\dagger}$ & 11280 & $^{\dagger}$ & 572 & 3538 \\
\midrule
\multicolumn{2}{@{}l}{\textbf{gmean} (14)} & 57.6 & 61.4 & 67.0 & & & & \\
\multicolumn{2}{@{}l}{\textbf{gmean} (16)} & 78.3 & 83.9 & & & \textbf{71.7} & & \\
\multicolumn{2}{@{}l}{\textbf{gmean} (12)} & 413.3 & & & 411.7 & & & \\
\bottomrule
\end{tabular}%
}
\end{table}

Table~\ref{tab:ablation} attributes \PORTER{}'s EPR consumption to its
mechanisms by disabling each in isolation, holding the initial layout
and every other parameter fixed; the per-action counts behind these
rows are tabulated in \appMix{} (the six omitted rows --- GHZ and
Graphstate at every scale --- form no packets, never trigger relief and
issue no telegates, so all configurations coincide there). Two columns
go further than a per-primitive ablation: ``Hop-wise'' replaces
\PORTER{}'s unit of commitment with \dSABRE{}'s while holding scoring,
plan set and relief tiers fixed, and ``\TeleData{}'' removes the
\TeleGate{} family outright, leaving state teleportation as the sole
inter-core action.

\textbf{Packets.} Disabling \TeleGatePacket{} makes every
shared-control remote CX run pay per gate, and the cost tracks the
packet count: $+42\%$ EPR on 64q~QFT (262 $\to$ 372, 45 packets),
$+40\%$ on 80q~AE (490 $\to$ 688, 264 packets), $+19\%$ on 64q~AE,
against a $+7\%$ overall gmean; the lone counter-move (80q~QFT,
$446 \to 435$) shows packet formation is a heuristic commitment that
can displace a slightly better teleport plan. Read alone, this column
is a clean win for the primitive. The next one overturns that reading.

\textbf{\TeleData{} only.} Removing the \TeleGate{} family shows two
things. First, the \TeleGate{} is what makes \PORTER{} \emph{terminate},
not merely what makes it cheap. It is the only remote action needing no
free destination slot, so it is always a legal plan for an
adjacent-remote CX. Without it, a front whose remote CXs are all
core-adjacent and whose destinations sit at or below $f_\mathrm{min}$
has an empty plan set, while the tail-relief BFS is offered only the
non-adjacent blocked candidates: ten of the twelve rows deadlock.
Feeding adjacent-blocked candidates into the same BFS (marked
$^{\ddagger}$, and inert whenever the \TeleGate{} is available) recovers
all but the two Random rows, which livelock --- and those two are the
most telegate-heavy in the suite ($228$ and $3{,}538$ single telegates).

Second, where \TeleData{}-only completes it is cheaper than the full
action set, increasingly so with scale: $+2\%$ at 25q, mixed at 64q,
$-33\%$ at 80q, and $8.5\%$ below Full over the sixteen instances it
completes. The circuits it wins are the packet-heaviest, so the ordering
there is $\text{\TeleData{} only} < \text{Full} < \text{No packets}$:
packets beat the single telegates they degrade into, but a pure
teleportation policy beats both. Section~\ref{sec:discussion} argues
this is a currency mismatch in the plan comparator rather than a verdict
on packets. It also answers a question \dSABRE{} leaves open. Its
teledata-only vocabulary was chosen by measurement --- an ungrouped
telegate serves one gate where teledata amortises many --- with grouping
untested. Grouping does close that per-EPR gap, and still does not
justify spending pairs that teleporting would have spent better.

\textbf{Plan commitment.} Replacing the commitment rule with
\dSABRE{}'s --- returning to the top of the loop after one round of
preparation SWAPs, re-classifying the front and re-scoring every plan
before the next move --- leaves EPR unchanged: $413.3$ against $411.7$
over the twelve rows here, and $148.8$ against $148.0$ over the
twenty-one instances we ran it on, with eight circuits better without
commitment, five worse, eight tied, and all twenty-one completing either
way. The two make the same number of routing moves ($128{,}222$ against
$128{,}565$ iterations), which is the key number: commitment does not
change what the compiler decides, only what each decision costs. It buys
$1.31\times$ compile time in geometric mean and $3.07\times$ on
80q~Random. Section~\ref{sec:discussion} notes what this does and does
not say about the plan abstraction beneath it.

\textbf{Capacity relief.} Disabling both tiers costs completion, not
just EPR pairs: 80q~AE, QFT, QNN and Random all fail, because with every
destination below $f_\mathrm{min}$ and no way to drain qubits toward
reservoirs the loop has no legal remote plan left. Where routing does
complete, dense circuits pay heavily ($+89\%$ on 25q~QFT, $+78\%$ on
25q~QNN, $+85\%$ on 64q~AE). Two instances are cheaper \emph{without}
relief --- 64q~QNN ($864 \to 600$) and marginally 25q~Random ($401 \to
395$) --- showing the opportunistic tier can overpay where congestion
would have cleared itself. Relief is thus the one mechanism here whose
removal costs completion: with both tiers on, no instance in this
evaluation reached the deadlock path of
Section~\ref{sec:sfc-progress}; with both off, four did. That is a
statement about these instances, not a guarantee. The initial layout,
which this table holds fixed, is assessed in
Sections~\ref{sec:shared_results} and~\ref{sec:heavyhex}.

\subsection{RQ4: Scale-Up on QFT and QPEexact to 360 Logical Qubits}
\label{sec:scaleup}

\emph{Question:} do the RQ1 and RQ2 answers hold as circuit width and core
count grow together? \emph{Held fixed:} \dSABRE{}'s own scalability
architectures and circuit files. \emph{Metric:} geometric-mean EPR and
complete compile time.


\begin{table}[t]
\centering
\caption{Scale-up suite: QFT and QPEexact at 100, 200, and 360 logical
  qubits, EPR pairs. Architectures follow \dSABRE{}'s own scalability
  protocol --- $5{\times}5$ cores throughout, core count $6 \to 12 \to
  20$ --- read from the baselines' own device files so all routers see
  identical physical qubit ids; our \dSABRE{} column reproduces the
  figures it publishes for exactly this series~\cite{lidsabre}
  ($201/246/766/896/1890/1938$) to within the one candidate layout our
  protocol adds. The last two columns are the fixed-layout control:
  \PORTER{} routing \dSABRE{}'s own \texttt{SabreLayout} candidates.
  ``$t$'' is one \dSABRE{} attempt (seconds); its protocol runs three
  passes over each of up to six candidates. \texttt{pytket-dqc} is
  omitted: its hypergraph partitioner does not finish in reasonable time
  at these widths.}
\label{tab:scaleup}
\setlength{\tabcolsep}{3.5pt}
\small
\resizebox{\columnwidth}{!}{%
\begin{tabular}{@{}ll rrr rr rr rr@{}}
\toprule
 & & \multicolumn{3}{c}{architecture} & \multicolumn{2}{c}{\dSABRE{}}
 & \multicolumn{2}{c}{\PORTER{}, own layout}
 & \multicolumn{2}{c}{\PORTER{}, \dSABRE{} layout} \\
\cmidrule(lr){3-5}\cmidrule(lr){6-7}\cmidrule(lr){8-9}\cmidrule(lr){10-11}
Circuit & CX & $n$ & $P$ & $K$ & EPR & $t$ & 1 pass & $+$R & 1 pass & $+$R \\
\midrule
QFT      &   3{,}420 & 100 & 150 &  6 &  200 &  1.7 &  198 &  339 &  240 &  273 \\
QPEexact &   3{,}579 & 100 & 150 &  6 &  246 &  2.0 &  200 &  262 &  200 &  283 \\
QFT      &   7{,}220 & 200 & 300 & 12 &  765 &  6.5 &  579 & 1128 & 1031 & 1018 \\
QPEexact &   7{,}579 & 200 & 300 & 12 &  896 &  7.8 &  528 &  979 &  651 &  946 \\
QFT      &  13{,}300 & 360 & 500 & 20 & 1890 & 29.9 & 1156 & 2822 & 1870 & 2546 \\
QPEexact &  13{,}979 & 360 & 500 & 20 & 1938 & 38.9 & 1224 & 2342 & 1444 & 2416 \\
\midrule
\textbf{gmean} (6) & & & & & \textbf{705.7} & & \textbf{507.7} & 930.3 & 665.7 & 877.9 \\
\bottomrule
\end{tabular}%
}
\end{table}

Every result so far is at 108 physical qubits or fewer.
Table~\ref{tab:scaleup} takes the two circuit families that exist at
100, 200 and 360 logical qubits onto \dSABRE{}'s scalability series:
$5{\times}5$ cores throughout, core count $6 \to 12 \to 20$, so
occupancy stays roughly constant while the core graph grows. It is a
monotone ladder in the quantity that matters --- inter-core traffic
scales with core count, and \dSABRE{}'s own cost rises $200 \to 765 \to
1{,}890$ on QFT.

\PORTER{} uses fewer EPR pairs than \dSABRE{} on all six instances, by a
widening margin: $0.90\times$ at 100 logical qubits, $0.67\times$ at
200, $0.62\times$ at 360, for $0.72\times$ overall. It does so at
$0.20\times$ \dSABRE{}'s complete compile cost ($8.4$\,s against
$41.1$\,s), because \dSABRE{} routes up to eighteen times here against
\PORTER{}'s three.

\textbf{The lookahead horizon carries that result, and nothing else
does.} Both mechanisms that close \PORTER{}'s gap below $108$ physical
qubits are inert or harmful at these widths. Backward refinement loses
on all six instances, by $83\%$. Handing \PORTER{} \dSABRE{}'s own
\texttt{SabreLayout} candidates --- the control that reaches parity at
25 and 64 qubits and wins on both heavy-hex core graphs --- makes it
\emph{worse} here, $0.94\times$ against $0.72\times$, and worst on the
largest instance, where \dSABRE{}'s layout costs $1{,}870$ pairs against
$1{,}156$ from \PORTER{}'s own. The mapping is not the binding
constraint at this scale; the horizon is. Every instance in this series
selects $L_\mathrm{tp}{=}128$, and at the configured $16$ \PORTER{}
trails \dSABRE{} by $1.45\times$ (Section~\ref{sec:horizon}).

This separates two deficits the smaller suites conflate. Up to $108$
physical qubits the initial layout limits \PORTER{}: control it, match
the budget, and the difference disappears. From $150$ upward the layout
is no longer the constraint and the horizon is. Selecting it per circuit
is a workaround, not a fix; a $K$-aware horizon remains open.

\subsection{RQ4: Heavy-Hex Cores and Non-Grid Core Graphs}
\label{sec:heavyhex}

\emph{Question:} how much of the foregoing depends on square cores wired
in a grid? \emph{Held fixed:} the 64-qubit circuit suite and the
physical-qubit count; the core shape and core graph change.
\emph{Metric:} geometric-mean EPR and win/loss counts.


\begin{table}[t]
\centering
\caption{The nine-circuit 64-qubit suite on four IBM 27-qubit
  heavy-hex cores (108 physical qubits),
  wired as a \emph{ring} (4 links, every core degree 2) and a
  \emph{star} (3 links, core~0 a hub every route must cross).
  Geometric-mean EPR; the right-hand pair is the fixed-layout control,
  \PORTER{} routing \dSABRE{}'s own \texttt{SabreLayout} candidates.
  \PORTER{} completes all eighteen instances, so every row is over the
  full nine; \PORTER{}+R, which does not select a lookahead horizon,
  still fails star Multiplier and is therefore ``---'' on that row.
  Per-circuit rows are in \appTables{}. \dSABRE{} reproduces its published numbers
  for these two architectures on all eighteen instances~\cite{lidsabre}.}
\label{tab:heavyhex}
\setlength{\tabcolsep}{4pt}
\small
\begin{tabular}{@{}l r rr rr@{}}
\toprule
 & & \multicolumn{2}{c}{own layout} & \multicolumn{2}{c}{\dSABRE{} layout} \\
\cmidrule(lr){3-4}\cmidrule(lr){5-6}
Core graph & \dSABRE{} & \PORTER{} & $+$R & \PORTER{} & $+$R \\
\midrule
ring (9)          &  85.2 &  86.1 & 82.3 &  82.9 & \textbf{81.4} \\
star (9)          & 101.1 & 122.3 &  --- &  98.5 & \textbf{93.5} \\
\midrule
star Multiplier   &   722 & \textbf{455} &  --- &   462 &   462 \\
\bottomrule
\end{tabular}
\end{table}


Every result so far uses square cores wired in a grid. To test how much
of \PORTER{}'s behaviour depends on that we adopt \dSABRE{}'s non-grid
experiment~\cite{lidsabre}: the same nine-circuit 64-qubit suite on four
IBM 27-qubit heavy-hex tiles. Two things change at once. The \emph{core}
becomes irregular --- degrees over $\{1,2,3\}$, six degree-1 leaves,
diameter 12 against a $4{\times}4$ grid's 6 --- and the \emph{core
graph} is either a ring, where every core has degree 2, or a star, where
core~0 is a hub. \PORTER{}'s hardware object is derived from the
architecture \dSABRE{}'s own builder produces, and a topology assertion
checks that both compilers see identical qubit ids, intra-core edges and
port endpoints before anything is routed. \dSABRE{} routes on its native
object and reproduces its published numbers on all eighteen instances.
Neither compiler is specialised for a heavy-hex tile.

\textbf{The core graph matters; the core shape does not.}
Table~\ref{tab:heavyhex} shows an irregular tile costs \PORTER{}
nothing: on the ring \PORTER{}+R leads \dSABRE{}
($82.3$ against $85.2$), winning six of nine and tying GHZ --- better
relative standing than on the grid suite. Doubling the intra-core
diameter and replacing degree-4 interiors with degree-1 leaves disturbs
neither compiler. The star does. Under its own protocol \PORTER{} trails
there, $122.3$ against $101.1$, winning two of nine; \PORTER{}+R, which
cannot route Multiplier at all, trails by $13.6\%$ on the eight it
completes.

\textbf{Star Multiplier is a livelock the horizon escapes.} At a fixed
$L_\mathrm{tp}{=}16$ this instance does not terminate, and it fails as a
local-SWAP livelock rather than the capacity deadlock of
Section~\ref{sec:sfc-progress}: retired gates freeze at $3{,}226$ of
$13{,}040$ and the EPR count at $538$ while intra-core SWAPs accumulate
at one per iteration, so the iteration guard rather than the progress
argument stops it. It is a property of the horizon, not of the circuit:
the same layout and the same router complete it at
$L_\mathrm{tp}{=}128$, at $455$ pairs against \dSABRE{}'s $722$. Since
horizon selection routes that candidate as a matter of course,
\PORTER{} has no unrouted instance anywhere in this evaluation ---
whereas \PORTER{}+R, which holds the horizon fixed, still fails it.

\textbf{The star gap is the layout, not the routing.} Given
\dSABRE{}'s layout, \PORTER{} completes all nine star circuits,
including Multiplier at $462$ pairs against \dSABRE{}'s $722$. With the
budget matched, that control reverses the result on \emph{both} core
graphs ($-4.4\%$ ring, $-7.5\%$ star), the larger margin on the topology
\PORTER{} loses under its own protocol. The cause is in the layout:
\PORTER{} reserves $f_\mathrm{min}$ in every core, treating cores as
interchangeable. They are on a grid and a ring, where every core has the
same degree. They are not on a star, where Multiplier opens with free
counts $(2,2,13,27)$ --- the hub every route must cross pinned at the
floor while a whole core sits empty. Raising the reservation uniformly
does not repair it (we tried $f_\mathrm{min}{\in}\{4,6\}$); what the
star needs is a reservation scaled by a core's position in the core
graph, which we have not implemented.

Taken with Section~\ref{sec:scaleup}, this bounds the layout reading
rather than extending it. Up to $108$ physical qubits --- across two
core shapes and three core graphs --- the initial layout separates the
two compilers and controlling it erases the difference. Above $150$ the
horizon takes over.

\subsection{The EPR/Time Operating Frontier}
\label{sec:frontier}
\label{sec:time}

\begin{table}[t]
\centering
\caption{Protocol trade-off over all 24 instances. ``routes'' is the number of
  routing passes; times are complete compile costs (geometric mean per
  instance). \PORTER{} routes each of three candidate lookahead
  horizons and keeps the best (Section~\ref{sec:horizon}); \PORTER{}+R adds backward
  refinement (Section~\ref{sec:sfc-refine}) at a fixed horizon;
  \PORTER{}+S adds the best-of-$\gamma$ layout search of this section
  on top of it; \PORTER{}+SR takes the better of horizon selection and
  that search, netting the one route they share. All are seed-free.}
\label{tab:protocols}
\small
\resizebox{\columnwidth}{!}{%
\begin{tabular}{@{}l r rr rr@{}}
\toprule
Protocol & routes & EPR & vs.\ \dSABRE{} & time (s) & vs.\ \dSABRE{} \\
\midrule
\dSABRE{}                 & 9--18 & 163.9 & $1.00$ & 2.26 & $1.00$ \\
\PORTER{}                 &     3 & 174.9 & $1.07$ & 1.32 & $\mathbf{0.58}$ \\
\PORTER{}+R               &     5 & 171.4 & $1.05$ & 1.67 & $0.74$ \\
\PORTER{}+S               & 9--13 & 163.2 & $1.00$ & 4.59 & $2.03$ \\
\PORTER{}+SR              & 11--15 & \textbf{158.2} & $\mathbf{0.97}$ & 5.65 & $2.50$ \\
\bottomrule
\end{tabular}%
}
\end{table}

Table~\ref{tab:protocols} places \PORTER{}'s four operating points
(Section~\ref{sec:opprotocols}) against \dSABRE{}, pooled over all
twenty-four instances. Selection and the $\gamma$-search look for
different things --- the horizon at a fixed layout against the layout at
a fixed horizon --- and neither dominates: the search gives the better
route on fourteen instances, selection on four, and they tie on six.
\PORTER{}+SR runs both.

The trade is time against EPR pairs and it is monotone in both.
\PORTER{} costs $1.07\times$ \dSABRE{}'s EPR at $0.58\times$ its compile
time; refinement reaches $1.05\times$ at $0.74\times$; the layout search
reaches parity ($163.2$ against $163.9$) at $2.03\times$; \PORTER{}+SR
passes it, at $0.97\times$ its EPR for $2.50\times$ its time.
\PORTER{}+SR leads \dSABRE{} on twelve instances against nine behind,
and by construction never loses to either schedule it is built from.
Increasing \PORTER{}'s deterministic search budget therefore closes and
reverses the own-protocol gap. We stop short of calling that gap a
search-budget gap: the added schedules change the initial layout as well
as the route, and the comparison equalises route counts rather than
wall-clock time. Where to sit on the curve depends on whether EPR pairs
or compile seconds are scarce.

Selecting the horizon has largely overtaken refinement as an operating
point. \PORTER{}+R is ahead by only $2.0\%$ in geometric mean for four
extra routes, and it is $12.0\%$ behind the selected route at 80 qubits,
$83\%$ behind on the scale-up series, and the configuration that fails
star~Multiplier. We report it because the fixed-layout comparison of
Section~\ref{sec:shared_results} needs a schedule that spends its extra
effort on routes rather than layouts, which \PORTER{}+S cannot be.

All times are \emph{complete} compile costs: layout generation plus
every route over every candidate, not the winning route alone. That is
nine to eighteen routes for \dSABRE{}, three for \PORTER{}, five for
\PORTER{}+R, and all three seeds for \TeleSABRE{}. \PORTER{}'s advantage
is in the count of routes, not their price. Per route it is the dearer
of the two on all nine 64-qubit circuits, at $0.27\times$ \dSABRE{}'s
speed in geometric mean: enumerating whole plans, collecting packets by
forward simulation and running a two-tier capacity search all cost more
per iteration than scoring a flat list of hops, and a deep horizon costs
more still. On that subset the complete costs are $1.98$\,s against
$1.86$\,s, so \dSABRE{} is cheaper; pooled over all twenty-four
\PORTER{} leads at $0.58\times$.

Over the five circuits \TeleSABRE{} completes, geometric-mean complete
cost is $5.95$\,s for \TeleSABRE{}, $2.67$\,s for \dSABRE{}, $1.55$\,s
for \PORTER{} and $1.97$\,s with refinement (per-circuit times in
\appTables{}). Neither \PORTER{} point holds that margin on the largest
circuit, where \PORTER{}+R costs $191$\,s against $144$\,s at
$23{,}381$ CX. The cheap operating point is still \PORTER{}, and at
$1.07\times$ \dSABRE{}'s EPR it gives up little to be so.
